% Optional inclusion fragment. Requires the B grammar of obstruction.tex.
% This theorem concerns standard term functions, not all definable relations.

\section{A stronger term invariant: only finite information from other inputs}
\label{sec:finite-control}

The coordinate bound leaves open whether the tied-seed language could
compute a bounded numerical comparison. This section closes that
particular route. The gate supplied by the current drafts does not
suffice to construct a numerical equality test, minimum, or truncated
subtraction. A stronger semantic invariant excludes all three.

\begin{definition}[One numerical input and finite control]
\label{fc:def}
A function $f:\mathbb N^r\to\mathbb N$, with $r\ge1$, has the finite
control property if there are an active coordinate $i$, finite nonempty
sets $A_j$ and unary maps $\alpha_j:\mathbb N\to A_j$ for $j\ne i$,
and unary functions $F_a:\mathbb N\to\mathbb N$ indexed by
$a\in\prod_{j\ne i}A_j$, such that
\[
 f(x_1,\ldots,x_r)
   =F_{(\alpha_j(x_j))_{j\ne i}}(x_i).
 \tag{FC}\label{fc:form}
\]
Thus the value may depend without a bound on one coordinate, while each
other coordinate supplies only a member of a fixed finite set. For
primitive-recursive functions, all displayed maps will be chosen
primitive recursive, using finite encodings of the sets $A_j$.
\end{definition}

A constant function satisfies the definition with any active coordinate.
Information about the active coordinate may also determine a finite
choice, since that choice can be absorbed into the unary function
$F_a$. The definition imposes no restriction at arity one.
Throughout the constructive argument, finite ranges may be replaced by
the specified finite codomains; no decision procedure for exact attained
ranges is assumed.

\begin{lemma}[Finite-valued functions and composition]
\label{fc:composition}
\begin{enumerate}
\item If $f$ has the finite control property and $\rho$ is unary with
      finite range, then $\rho\circ f$ factors through finite-valued
      unary maps of \emph{all} input coordinates.
\item The functions with the finite control property contain constants
      and projections and are closed under composition.
\end{enumerate}
Both assertions preserve primitive recursiveness of the displayed maps.
\end{lemma}

\begin{proof}
For the first assertion, use~\eqref{fc:form} and attach to the active
coordinate the finite feature
\[
                 \beta_i(x)=\bigl(\rho(F_a(x))\bigr)_a.
\]
There are finitely many $a$, and each entry lies in the finite range of
$\rho$. Together with the existing $\alpha_j(x_j)$, this determines
$\rho(f(\vec x))$.

For composition write $f(g_1(\vec x),\ldots,g_r(\vec x))$, where $f$
has active input $i$ and $g_i$ has active input $k$. Every other $g_j$
enters $f$ only through its finite-valued feature $\alpha_j\circ g_j$.
By the first assertion these features are determined by finite unary
features of all coordinates of $\vec x$. Combine their features on
coordinates other than $k$ with the finite controls of $g_i$. Once those
combined controls are fixed, the whole expression is one unary function
of $x_k$: features depending on $x_k$ are evaluated inside that unary
function. There are only finitely many combined controls, giving the
required form. All vector formation, finite selection, and composition
here are primitive recursive when the original maps are.
\end{proof}

\begin{theorem}[Finite control for every tied-seed term]
\label{fc:B-theorem}
Every standard numerical function of positive arity represented by a
$B$-term in the grammar of
Section~\ref{sec:finite-dimensional-obstruction} has the finite control
property. Nullary terms are constants and satisfy the same conclusion
after an unused input is added. A finite representation of the form
\eqref{fc:form}, with primitive-recursive feature maps and profiles, can
be constructed by induction on its syntax.
\end{theorem}

\begin{proof}
Variables and zero are covered by Lemma~\ref{fc:composition}. It is
enough to show that the operation $B_h(n,\vec p)$ preserves the property
when its fixed update $h(u,v,\vec p)$ has it. The actual counter term
is then substituted using the composition lemma.

\paragraph{The active input of $h$ is its previous state.}
Write
\[
                 h(u,v,\vec p)=F_{\alpha(u),\gamma(\vec p)}(v),
\]
where $\gamma(\vec p)$ is a tuple of finite unary features, one for
each passive coordinate. For each of its finitely many values $c$, define
one unary sequence by
\[
 R_c(0)=1+F_{\alpha(0),c}(0),\qquad
 R_c(m+1)=F_{\alpha(m),c}(R_c(m)).
\]
These are primitive-recursive unary functions. Directly from the
standard defining equations,
\[
                      B_h(n,\vec p)=R_{\gamma(\vec p)}(n).
\]
This has active coordinate $n$ and finite controls on the passive inputs.

\paragraph{The active input of $h$ is its old counter.}
Write
\[
                 h(u,v,\vec p)=F_{\beta(v),\gamma(\vec p)}(u).
\]
For each finite passive profile $c$, define
\[
 R_c(0)=1+F_{\beta(0),c}(0),\qquad
 R_c(m+1)=F_{\beta(R_c(m)),c}(m).
\]
Again $B_h(n,\vec p)=R_{\gamma(\vec p)}(n)$, of the required form.

\paragraph{The active input of $h$ is a passive parameter $p$.}
Write the remaining passive parameters as $\vec q$ and put
\[
                 h(u,v,p,\vec q)
                    =F_{\alpha(u),\beta(v),\gamma(\vec q)}(p).
\]
The ranges of $\alpha$, $\beta$, and $\gamma$ are finite. Form one
finite catalogue of unary functions $V_j(p)$ containing all profiles
$F_{a,b,c}(p)$ and also every possible initial profile
\[
                    V_{0,c}(p)=1+F_{\alpha(0),\beta(0),c}(p).
\]
Give these occurrences separate indices if their extensional functions
happen to coincide. Denote the finite index set by $J$. Every state of
the recursion is the value at $p$ of one function in this catalogue.

The only observation of a previous state used by the update is its
$\beta$-value. Record all such observations at once in the finite unary
feature
\[
                         \kappa(p)=\bigl(\beta(V_j(p))\bigr)_{j\in J}.
\]
For each passive control value $c$ and each possible finite vector $k$,
define a finite-state process. Its initial state is the catalogue index
of $V_{0,c}$. At counter $m$, if its present index is $j$, its next index
is the catalogue occurrence of
\[
                              F_{\alpha(m),k_j,c}.
\]
Let $\theta_{c,k}(n)\in J$ be the index after $n$ transitions. This is
a finite-valued primitive-recursive unary function of $n$. It is
defined for all possible vectors $k$, including ones not realized by
any $p$; only the realized vectors will be used in the equality below.
Pack the finitely many processes into one finite-valued unary feature
\[
                              \eta(n)=\bigl(\theta_{c,k}(n)\bigr)_{c,k}.
\]
Induction on the standard counter in the metatheory gives
\[
 B_h(n,p,\vec q)
   =V_{\eta(n)_{\gamma(\vec q),\kappa(p)}}(p).
 \tag{Catalogue evaluation}\label{fc:catalogue}
\]
The initial case selects $V_{0,c}$; the transition uses exactly the
recorded observation $\beta(V_j(p))=\kappa(p)_j$.

For each fixed $e$ in the finite product codomain chosen for $\eta$
and each fixed $c$,
define the unary profile
\[
                              H_{e,c}(p)=V_{e_{c,\kappa(p)}}(p).
\]
Equation~\eqref{fc:catalogue} becomes
\[
                       B_h(n,p,\vec q)=H_{\eta(n),\gamma(\vec q)}(p).
\]
This has active coordinate $p$, with finite controls on $n$ and each
remaining passive input. It proves the last case. The construction is
finite and primitive recursive at every stage, completing the syntax
induction.
\end{proof}

The numerical bounds of Lemma~\ref{obs:coordinate-bound} are also an
immediate consequence of this stronger representation: take the maximum
of the finitely many unary profiles and then its running maximum. Its
new content is that the other inputs supply only finitely many distinct
profiles, rather than arbitrarily complicated bounded information.

\begin{corollary}[No numerical equality test, minimum, or variable countdown]
\label{fc:excluded}
None of the following binary functions is represented by a $B$-term:
\[
 \operatorname{eq}(a,b)=
 \begin{cases}1,&a=b,\\0,&a\ne b,\end{cases}
 \qquad
 \min(a,b),
 \qquad
 a\mathbin{\dotminus}b=\max(a-b,0).
\]
The same exclusion holds for the complementary numerical equality test.
\end{corollary}

\begin{proof}
For a binary function with active first coordinate, fixing the second
coordinate produces at most finitely many unary functions of the first.
If the second coordinate is active, the reverse collection of unary
profiles is finite. Each displayed function has infinitely many
distinct profiles in both directions.

For numerical equality, fixing one input at $b$ gives the characteristic
function of the singleton $\{b\}$ in the other input. For minimum,
the functions $x\mapsto\min(x,b)$ are distinct for distinct $b$, as is
seen by evaluating at an $x$ at least as large as both candidate values.
For truncated subtraction, distinct fixed second arguments give distinct
functions by evaluating the first argument above both; distinct fixed
first arguments give distinct functions already at second argument zero.
The finite-profile conclusion of Theorem~\ref{fc:B-theorem} is therefore
impossible in either orientation.
\end{proof}

\subsection{The exact consequence for the proposed relational route}

The current $B$ grammar does define ordinary predecessor and a one-sided
gate. Write
\[
 D(n,a)=B(n;u,v.a),\quad W(n)=B(n;u,v.u),\quad
 P(n)=W(D(n,n)),\quad Z(n)=D(n,0).
\]
Their computation equations give
$D(0,a)=Sa$, $D(Sn,a)=a$, $P0=0$, $P(Sn)=n$, $Z0=1$, and $Z(Sn)=0$.
With equation induction, the gate
\[
             G(n,a)=P\bigl(P(B(a;u,v.D(Z(n),v)))\bigr)
\]
satisfies $G(0,a)=0$ and $G(Sn,a)=a$. At zero its internal update is
the identity and its state is constantly $1$; at a positive control its
update is successor and its state is $S^2a$.

For completeness, both gate identities are internal consequences of the
$B$ computation schemas, ordinary equality reasoning, and literal-equation
induction. For the zero control, induct directly on
\[
                 B(a;u,v.D(Z(0),v))=1.
\]
Its base is $S(D(1,0))=S0=1$; its successor computation applies
$D(1,\cdot)$, which repeats the previous state. For the positive control
$Sq$, induct directly on
\[
                 B(a;u,v.D(Z(Sq),v))=S^2a.
\]
Its base is $S(D(0,0))=S(S0)=2$; its successor computation applies
$D(0,\cdot)$, hence successor, to the previous state. Keeping these
actual update bodies throughout avoids invoking an unmentioned
extensionality principle for recursion descriptions. The displayed
equations for $P$ then remove the two outer predecessors.
All control parameters are fixed during these inductions; no induction
on a disjunction or a compound formula is used.

\begin{lemma}[A positive boundary: arbitrary zero-start counter recursion]
\label{fc:zero-start-lemma}
In any such equational presentation, the gate gives a general zero-start
counter recursor as a macro:
\[
 L(n;u,v.h)
   :=B\bigl(Sn;r,s.G(r,h[Pr/u,s/v])\bigr),
 \tag{Zero-start macro}\label{fc:zero-start}
\]
with all new bound variables fresh. It satisfies
\[
 L(0;u,v.h)=0,\qquad
 L(Sn;u,v.h)=h[n/u,L(n;u,v.h)/v].
\]
\end{lemma}

\begin{proof}
Let $H(r,s)=G(r,h[Pr/u,s/v])$ be the fixed internal update. The gate
equations and $P(Sr)=r$ give
\[
               H(0,s)=0,\qquad H(Sr,s)=h[r/u,s/v].
\]
Thus $B_H(0)=S(H(0,0))=1$ and
$B_H(1)=H(0,B_H(0))=0$. The latter identity is the base equation for
$L(n)=B_H(Sn)$. At a successor,
\[
 L(Sn)=B_H(S(Sn))=H(Sn,B_H(Sn))=h[n/u,L(n)/v].
\]
No further induction is used after the gate identities have been proved.
\end{proof}

These statements do not conflict with Theorem~\ref{fc:B-theorem}:
their macros remain $B$-terms. In particular, fixed-modulus counting
and fixed finite-state controls do not yield a numerical equality
test for two unbounded independent inputs.

The theorem does not settle whether every unary primitive-recursive
function is represented, because its conclusion imposes no restriction
on unary functions. It also does not rule out formula-defined comparison
relations: the source already has logical equality $a=b$, which must
not be confused with a term returning its numerical characteristic.
This term theorem alone leaves quantified operation graphs outside its
scope. Section~\ref{sec:relational-obstruction} adds a proof-theoretic
argument and excludes arbitrary finite-dimensional relational
bi-interpretations for the ordinary tied-seed theory with its stated
existential-free axiomatization.
